\section{The retained inputs and the two endpoint topologies}
\label{sec:v91-inputs}
All densities below are those of the original exact event. For
$q>1$, an interval $J$, and a measurable function $f$, write
\begin{equation}\label{eq:v91-norm}
 \mathcal N_q(f;J)=\sup_{E\subset J,\,0<|E|<\infty}
                  |E|^{-1+1/q}\int_E|f(t)|\,dt.
\end{equation}
This is the measurable-set norm for weak $L^q$, not the strong
$L^q$ norm. Set $Q=32/31$, $\gamma=1/32$ and $\alpha=129/128$.
The word ``critical'' below refers to this finite exponent, never
to the essential-supremum target \eqref{eq:v76-target}.

\begin{inputstatement}[Original source comparison]
\label{input:v91-comparison}
The physical stopped flow has $XF=1$, with $F=L$, and preserves
the endpoint-momentum angle. Its maximal regular segments have total
residence $\ell=\tau^-+\tau^+$, gradient magnitude $z=|\nabla_qF|$,
and physical angular density $\omega/(4\pi c_*\chi)$. Nongrazing
outgoing faces have the same one-sided density, their incident
components are retained, and their complete first-exit assignment is
\[
 \mathcal F_h(t)=\int_{\Gamma_{\rm ng}^+}
 \one_{\{0<b-t\le h,\ b-t<\tau^-(y)\}}\,d\zeta_+(y),
 \qquad \overline{\mathcal F}_r=\frac1r\int_0^r\mathcal F_h\,dh.
\]
Here $b$ denotes the outgoing roof, and $d\zeta_+$ is the physical
outgoing flux. Put $I_m=1+\sum_{i=1}^{m-1}q_0^{d_i}/c_i$ and
\[
 v^{\eta,a}(t)\,dt=F_\#(I_mz^{-\eta}\one_{\{0<z<a\}}\one_\Pi\nu_R^*),
 \quad v^\eta(t)\,dt=F_\#(I_mz^{-\eta}\one_\Pi\nu_R^*).
\]
For $\eta=1/4$, let $v=v^{1/4}$ and
$\mathsf P_rv(t)=\int_0^r s^{-3/8}v(t+s)\,ds$. The complete
source satisfies
\begin{align}
 b^\varepsilon(t)\le{}&2A_rb^\varepsilon(t)
 +C\overline{\mathcal F}_r(t)+C\mathsf P_rv(t)\notag\\
 &+C_\Theta(1+\varepsilon^{-1})\sqrt r\,A_rp(t),
 \qquad A_rf(t)=\frac1{2r}\int_{t-r}^{t+r}f(u)\,du.
 \label{eq:v91-source-comparison}
\end{align}
This holds for each fixed $r>0$ at almost every roof, without a
positive incidence, gradient or residence cutoff.
\end{inputstatement}
This is the retained physical source-shift inequality combined with
the two-sided optical payment of the actual first-exit surface.
True exits and initial weights are included before the comparison;
the surface is not a volume collar.

\begin{inputstatement}[Local estimates and positive remainder]
\label{input:v91-local}
For every previously fixed positive interval length and every fixed
$\varepsilon>0$,
\begin{align}
 \limsup_m\sup_{R,n,k}m^2\|A_r b^\varepsilon\|_\infty
                                      &\le C\varepsilon^{1/16},\notag\\
 \limsup_m\sup_{R,n,k}m^2\|A_rp\|_\infty&\le C,\notag\\
 \limsup_m\sup_{R,n,k}\sup_{|J|=L}m^2\int_Jv
                                      &\le CL,\notag\\
 \limsup_m\sup_{R,n,k}\sup_{|J|=L}m^2\int_Jv^{1/4,a}
                                      &\le Ca^{\alpha-1/4}L,\notag\\
 \limsup_m\sup_{R,n,k}\sup_{|J|=L}m^2
   \int_J\!\int\one_{\{\ell<\delta\}}\,d\mu_t\,dt
                                      &\le C\delta^\gamma L.
 \label{eq:v91-local-input}
\end{align}
The leading constants are independent of the previously fixed widths.
The retained critical surface theorem, evaluated at window length
$r$, gives
\begin{equation}\label{eq:v91-matched-input}
 \limsup_m\sup_{R,n,k}\sup_{|I|=r}
       \mathcal N_Q(m^2\overline{\mathcal F}_r;I)\le Cr.
\end{equation}
Its underlying estimate is
$C\{L^{31/32}r^{1/32}+(L+4r)(1+r^{19/32})\}$ on windows
of length $L$. It sums the complete capped residence tail at finite
collision count. No logarithmic norm occurs in this bound.

For sufficiently large fixed $B$, with $\varepsilon(B)=A_0B^{-1/12}$,
\begin{equation}\label{eq:v91-positive-remainder}
 m^2p-\mathcal L=m^2b^{\varepsilon(B)}+e_B,
 \qquad\limsup_m\sup_{R,n,k}\esssup_{\mathcal C(M)}|e_B|
                                        \le C_MB^{-1/192}.
\end{equation}
The finite-peak reference satisfies
$|\mathcal L|\le D_M$ on $\mathcal C(M)$, uniformly in all labels,
radii and collision counts. Critical points have zero source measure.
\end{inputstatement}
The source comparison and moment inputs are restated in the retained
article's geometric-input section. The matched estimate is precisely
the conclusion of its critical-surface corollary, not a new surface
spectral hypothesis. The new proofs below add only finite interval
gluing, positive source comparison and real-variable maximal estimates.
The full four-record spectral chain remains an inherited input.

For each $m,R,n,k$, the Borel maximal-segment construction fixes common
roof kernels outside a null set and common roof--time kernels outside
a joint null set. Every positive restriction uses those kernels.
The parameter convention is $\sup_R\esssup_t$, not one common null
set for all radii. In each inherited local theorem the auxiliary
band is fixed before $m\to\infty$ and enlarged afterwards. In the
new scale-gluing argument $r$ and the number of its subintervals are
also fixed before that limit; $B$ in
\eqref{eq:v91-positive-remainder} remains outermost.
